\TieudeT{PHONG TỎA VẬT LÍ 12 - VẬT LÍ NHIỆT}
\subsubsection*{PHẦN I. Thí sinh trả lời câu hỏi từ câu 1 đến câu 18. Mỗi câu hỏi thí sinh chỉ chọn một phương án}
\setcounter{ex}{0}
\begin{ex} %Câu 1
	"Độ không tuyệt đối" là nhiệt độ ứng với
	\choice
	{\True $0\ K$}
	{$0^\circ C$}
	{$273^\circ C$}
	{$273\ K$}
	\loigiai{}
\end{ex}

\begin{ex} %Câu 2
	Quá trình chất chuyển từ thể rắn sang thể lỏng là
	\choice
	{ngưng tụ}
	{hóa hơi}
	{đông đặc}
	{\True nóng chảy}
	\loigiai{}
\end{ex}

\begin{ex} %Câu 3
	Chất rắn có tính chất nào sau đây?
	\choice
	{Có thể nén được dễ dàng}
	{Không có thể tích riêng}
	{\True Có hình dạng riêng xác định}
	{Không có hình dạng riêng xác định}
	\loigiai{}
\end{ex}

\begin{ex} %Câu 4
	Trong thí nghiệm Brown, các hạt phấn hoa chuyển động là do
	\choice
	{gió thổi làm hạt phấn hoa chuyển động}
	{các hạt phấn hoa tự chuyển động}
	{giữa các hạt phấn hoa có khoảng cách}
	{\True các phân tử nước chuyển động va chạm vào}
	\loigiai{}
\end{ex}

\begin{ex} %Câu 5
	Gọi $x, y, z$ lần lượt là khoảng cách trung bình giữa các phân tử của một chất (không phải là nước) ở thể rắn, lỏng và khí. Phát biểu đúng là
	\choice
	{$z<y<x$}
	{$x<z<y$}
	{$y<x<z$}
	{\True $x<y<z$}
	\loigiai{}
\end{ex}

\begin{ex} %Câu 6
	Thả vào chậu nước có nhiệt độ $t_1$ một thỏi đồng được đun nóng đến nhiệt độ $t_2$ ($t_2>t_1$). Sau khi cân bằng nhiệt, cả hai có nhiệt độ $t$, khi đó
	\choice
	{$t>t_2>t_1$}
	{\True $t_2>t>t_1$}
	{$t_1>t>t_2$}
	{Không thể so sánh được}
	\loigiai{}
\end{ex}

\begin{ex} %Câu 7
	Người ta dùng một lực nén từ từ pit-tông của một xilanh đi xuống để thay đổi nhiệt năng của khí trong xilanh. Họ đã thay đổi nhiệt năng của khí bằng cách
	\choice
	{truyền nhiệt}
	{\True thực hiện công}
	{không xác định được}
	{vùa thực hiện công vừa truyền nhiệt}
	\loigiai{}
\end{ex}

\begin{ex} %Câu 8
	Nguyên lý truyền nhiệt là
	\choice
	{nhiệt tự truyền từ vật có nhiệt độ thấp hơn sang vật có nhiệt độ cao hơn}
	{\True nhiệt tự truyền từ vật có nhiệt độ cao hơn sang vật có nhiệt độ thấp hơn}
	{nhiệt truyền từ vật có nhiệt dung riêng cao hơn sang vật có nhiệt dung riêng thấp hơn}
	{nhiệt truyền từ vật có nhiệt dung riêng thấp hơn sang vật có nhiệt dung riêng cao hơn}
	\loigiai{}
\end{ex}

\begin{ex} %Câu 9
	Phát biểu nào sau đây là đúng?
	\choice
	{\True Chất khí không có hình dạng và thể tích xác định}
	{Chất lỏng không có thể tích riêng xác định}
	{Lực tương tác giữa các nguyên tử, phân tử trong chất khí là rất mạnh}
	{Trong chất lỏng các nguyên tử, phân tử dao động quanh vị trí cân bằng cố định}
	\loigiai{}
\end{ex}

\begin{ex} %Câu 10
	Khi nhiệt độ của vật tăng lên thì
	\choice
	{\True động năng trung bình của các phân tử cấu tạo nên vật tăng}
	{thế năng của các phân tử cấu tạo nên vật tăng}
	{động năng trung bình của các phân tử cấu tạo nên vật giảm}
	{nội năng của vật giảm}
	\loigiai{}
\end{ex}

\begin{ex} %Câu 11
	Nhiệt độ không khí vào một ngày thời tiết đẹp là $27^\circ C$. Trong nhiệt giai Kelvin, nhiệt độ này là
	\choice
	{$264\ K$}
	{\True $300\ K$}
	{$200\ K$}
	{$246\ K$}
	\loigiai{}
\end{ex}

\begin{ex} %Câu 12
	Một động cơ nhiệt có hiệu suất là $36\%$, động cơ nhận nhiệt lượng $9 \times 10^4\ J$. Công cơ học mà động cơ thực hiện là
	\choice
	{\True $32,4\ kJ$}
	{$250\ kJ$}
	{$57,6\ kJ$}
	{$4\ kJ$}
	\loigiai{}
\end{ex}

\begin{ex} %Câu 13
	Một nhiệt lượng $6900\ J$ được dùng để đốt nóng $500\ g$ sắt.  Biết $c = 460\ J/kg.K$. Độ tăng nhiệt độ của sắt là
	\choice
	{$20^\circ C$}
	{$25^\circ C$}
	{\True $30^\circ C$}
	{$35^\circ C$}
	\loigiai{}
\end{ex}

\begin{ex} %Câu 14
	Người ta truyền cho khí trong một xilanh nhiệt lượng $250\ J$. Khí nở ra thực hiện công $100\ J$ đẩy pit-tông dịch chuyển. Độ biến thiên nội năng của khí là
	\choice
	{$\Delta U=-150\ J$}
	{\True $\Delta U=150\ J$}
	{$\Delta U=-350\ J$}
	{$\Delta U=350\ J$}
	\loigiai{}
\end{ex}

\begin{ex} %Câu 15
	Pha $100\ g$ nước ở $100^\circ C$ vào $100\ g$ nước ở $40^\circ C$. Nhiệt độ cuối cùng của hỗn hợp nước là
	\choice
	{$30^\circ C$}
	{\True $70^\circ C$}
	{$50^\circ C$}
	{$60^\circ C$}
	\loigiai{}
\end{ex}

\begin{ex} %Câu 16
	Giả sử cung cấp cho hệ một công là $200\ J$ nhưng nhiệt lượng thất thoát ra ngoài là $120\ J$. Độ biến thiên nội năng của hệ là
	\choice
	{Giảm $320\ J$}
	{\True Tăng $80\ J$}
	{Giảm $80\ J$}
	{Tăng $320\ J$}
	\loigiai{}
\end{ex}

\begin{ex} %Câu 17
	Thả $20\ g$ nước đá ở $-20^\circ C$ vào bình cách nhiệt chứa $100\ g$ nước ở $20^\circ C$. Biết nhiệt dung riêng của nước đá, nhiệt nóng chảy của nước đá và nhiệt dung riêng của nước lần lượt là $c_{\text{đá}} = 2,1\ J/(g.K)$, $\lambda = 330\ J/g$, $c_{\text{nước}} = 4,2\ J/(g.K)$. Bỏ qua sự trao đổi nhiệt với bình. Nhiệt độ sau cùng của nước trong bình là
	\choice
	{$2,1^\circ C$}
	{\True $1,9^\circ C$}
	{$0^\circ C$}
	{$2,3^\circ C$}
	\loigiai{
	\begin{itemize}
		\item Nhiệt lượng nước tỏa ra khi giảm về $0^\circ C$: $Q_n = m_n.c_n.\Delta t_n = 0,1.4200.20 = 8400\ J$.
		\item Nhiệt lượng nước đá thu vào để tăng nhiệt độ và nóng chảy hoàn toàn: $Q = m_\text{đ}(\lambda + c_\text{đ}.\Delta t_\text{đ}) = 0,02 (3,3.10^5+2100.20) = 7440\ J$.
		\item Nhận thấy $Q_n>Q \Rightarrow$ nước đá sau khi nóng chảy hoàn toàn sẽ tăng nhiệt độ.
		\item Phương trình cân bằng nhiệt: $Q_n = Q_\text{đ} \Rightarrow 0,1.4200.(20-t_{cb}) = 7440 + 0,02.4200.t_{cb} \Rightarrow t_{cb} \approx 1,91^\circ C$.
	\end{itemize}
	}
\end{ex}

\begin{ex} %Câu 18
	Đổ $m\ (g)$ nước ở $43^\circ C$ vào bình cách nhiệt chứa $m\ (g)$ nước đá ở $-20^\circ C$. Khi cân bằng nhiệt, thu được hỗn hợp nước và nước đá ở $0^\circ C$. Biết nhiệt dung riêng của nước, nhiệt dung riêng của nước đá và nhiệt nóng chảy riêng của nước đá lần lượt là $c_{\text{nước}} = 4,2\ J/(g.K)$, $c_{\text{đá}} = 2,1\ J/(g.K)$, $\lambda = 330\ J/g$. Khối lượng nước đá còn lại bằng bao nhiêu phần trăm ban đầu?
	\choice
	{$26\%$}
	{$42\%$}
	{\True $58\%$}
	{$74\%$}
	\loigiai{
	Chuẩn hóa $m=1$
	\begin{itemize}
		\item $Q_{thu} = Q_{\text{tỏa}} \Rightarrow m_\text{đ}.c_\text{đ}.\Delta t_\text{đ} + (1-\Delta m).\lambda = m_n.c_n.\Delta t_n \Rightarrow 1.2,1.20 + (1-\Delta m).330 = 1.4,2.43 \Rightarrow \Delta m = 0,58 = 58\%$.
	\end{itemize}
	}
\end{ex}
\subsubsection*{PHẦN 2. Thí sinh trả lời từ câu 1 đến câu 4. Trong mỗi ý a), b), c), d) ở mỗi câu, thí sinh chọn đúng hoặc sai.}
\setcounter{ex}{0}
\begin{ex}
	Hình vẽ bên là đường biểu diễn sự thay đổi nhiệt độ theo thời gian của một chất ban đầu ở thể rắn.
	\immini{\choiceTFt
	{Nhiệt độ sôi của chất này là $80^\circ C$}
	{\True Trong $4$ phút đầu, nhiệt độ của chất này tăng dần}
	{Để đưa nhiệt độ chất này từ $50^\circ C$ tới $80^\circ C$ cần thời gian $30$ phút}
	{Trong khoảng thời gian mà nhiệt độ là $80^\circ C$ thì chất này tồn tại ở thể lỏng và thể hơi}
	}{\begin{tikzpicture}[thick,>=latex']
			\draw [->] (0,0) -- (6,0) node [above] {Thời gian (phút)};
			\draw [->] (0,0) -- (0,3.5) node [right] {Nhiệt độ $(^\circ C)$};
			\draw [dotted] (0,0) grid [step = 0.5] (5.5,3);
			\draw [line width=1.5pt] (0,0.5) -- (2,2) -- (3.5,2) -- (5.5,3);
			\foreach \x in {50,60,...,100}{
				\draw node [left] at (0,\x/20-2) {$\x$};
			}
			\foreach \x in {0,1,...,11}{
				\draw node [below] at (\x/2,0) {$\x$};
			}
	\end{tikzpicture}}
	
	\loigiai{
		
		\begin{itemchoice}
			\itemch Nhiệt độ nóng chảy là $80^\circ C$.
			\itemch
			\itemch 4 phút.
			\itemch Có cả thể rắn.
		\end{itemchoice}
	}
\end{ex}

\begin{ex}
	Một miếng đồng có khối lượng $400\ gam$ được đun nóng từ $30^\circ C$ đến $60^\circ C$ rồi bỏ miếng đồng đó vào một chất lỏng có khối lượng $m =2\ kg$ thì nhiệt độ của chất lỏng tăng thêm $0,54^\circ C$, còn nhiệt độ của miếng đồng trở lại $30^\circ C$. Nhiệt dung riêng của đồng là $c_{\text{đ}}=378\ J/kg.K$ và nhiệt dung riêng của rượu là $c_r=2500\ J/kg.K$
	\choiceTFt
	{Nhiệt lượng miếng đồng thu vào để từ $30^\circ C$ nóng lên $60^\circ C$ là $4560\ J$}
	{Nhiệt độ ban đầu của chất lỏng là $30,54^\circ C$}
	{\True Nhiệt dung riêng của chất lỏng đó là $4200\ J/kg.K$}
	{Chất lỏng đó là rượu}
	\loigiai{
		\begin{itemchoice}
			\itemch $Q = m.c.\Delta t = 0,4.378.30 = 4536\ J$.
			\itemch Nhiệt độ ban đầu của chất lỏng là $30-0,54 = 29,46^\circ C$.
			\itemch $m_{\text{đ}}.c_{\text{đ}}.\Delta t_{\text{đ}} = m.c.\Delta t \Rightarrow c = \dfrac{m_{\text{đ}}.c_{\text{đ}}.\Delta t_{\text{đ}}}{m.\Delta t} = \dfrac{4536}{2.0,54} = 4200\ J/kg.K$.
			\itemch 
		\end{itemchoice}
	}
\end{ex}

\begin{ex}
	\immini[thm]{Một động cơ nhiệt hoạt động giữa hai nguồn nhiệt có nhiệt độ $T_1$ và $T_2$ (như hình vẽ). Động cơ nhận nhiệt lượng $Q_1=4000\ J$ từ nguồn nóng trong mỗi giây, công suất của động cơ là $1,2\ kW$.
		\choiceTFt
		{\True $T_1>T_2$}
		{Công của động cơ trong mỗi giây là $1,2\ J$}
		{\True Nhiệt lượng tỏa ra cho nguồn lạnh trong mỗi giây là $2800\ J$}
		{\True Hiệu suất của động cơ là $30\%$}}{\begin{tikzpicture}[thick,>=latex']
			%\draw [dotted] (-3,-3) grid (3,3);
			\draw [fill = red] (-1.5,2) rectangle (1.5,1.3);
			\draw [fill = cyan] (-1.5,-1.3) rectangle (1.5,-2);
			\fill [ball color= gray!40] (0,0) circle (0.8cm);
			\fill [top color = red, bottom color = cyan!50] (0,1) --(-.5,1.4) -- (-0.5,-1.3) -- (0,-1.7) --(0.5,-1.3) -- (0.5,-0.2) -- (1.5,-0.2) -- (1.8,0) -- (1.5,0.2) -- (0.5,0.2) -- (0.5,1.4) -- cycle;
			\draw (0,1) --(-.5,1.4) -- (-0.5,-1.3) -- (0,-1.7) --(0.5,-1.3) -- (0.5,-0.2) -- (1.5,-0.2) -- (1.8,0) -- (1.5,0.2) -- (0.5,0.2) -- (0.5,1.4) -- cycle;
			\draw 
			node [above] at (0,2) {Nguồn nóng $T_1$}
			node [below] at (0,-2) {Nguồn lạnh $T_2$}
			node [right] at (0.5,1) {$Q_1$}
			node [right] at (0.5,-1) {$Q_2$}
			node [right] at (1.8,0) {$A$}
			(-0.8,0) -- (-2,0) node [above, text width=2.5cm, align=center] {Tác nhân và cơ cấu động cơ nhiệt};
			
	\end{tikzpicture}}
	\loigiai{
		\begin{itemchoice}
			\itemch 
			\itemch 
			\itemch 
			\itemch 
		\end{itemchoice}
	}
\end{ex}

\begin{ex}
	\immini[thm]{Một bếp điện có công suất $P=1\ kW$. Tại thời điểm $\tau =0$ bắt đầu đun một lượng nước có nhiệt độ ban đầu là $20^\circ C$. Đồ thị hình bên biểu diễn sự phụ thuộc nhiệt độ của nước vào thời gian đun. Biết từ thời điểm $5$ phút đến thời điểm $8$ phút bếp bị mất điện, nhiệt lượng nước tỏa ra môi trường ngoài tỉ lệ thuận với thời gian và nhiệt dung riêng của nước là $c_n=4200\ J/kg.K$.}{\begin{tikzpicture}[thick,>=latex']
			%\draw [dotted] (-3,-3) grid (3,3);
			\draw [->] (0,0) -- (0,5.5) node [right] {$t(^\circ C)$};
			\draw [->] (0,0) -- (5.5,0) node [right] {$\tau$(phút)};
			\draw [line width=1.5pt] (0,1) node [left] {$20$} -- (1.25,2.25) -- (2,2) -- (5,5);
			\draw [dashed]
			(0,2.25) node [left] {$45$} -- (1.25,2.25) -- (1.25,0) node [below] {$5$}
			(0,2) node [left] {$40$} -- (2,2)-- (2,0) node [below] {$8$}
			(0,5) node [left] {$100$} -- (5,5) -- (5,0) node [below] {$\tau_3$};
	\end{tikzpicture}}
	
	\choiceTFt
	{Khi mất điện nhiệt độ của nước không đổi}
	{Công suất tỏa nhiệt ra môi trường bên ngoài của nước là $200\ W$}
	{Khối lượng nước ban đầu là $3,4\ kg$}
	{$\tau_3=16$ (phút)}
	\loigiai{
		\begin{itemchoice}
			\itemch Khi mất điện nhiệt độ nước giảm
			\itemch
			$\begin{cases}
				\left(\mathcal{P} - \mathcal{P}_{hp} \right).5.60  &= m.c.(45-20)\\
				\mathcal{P}_{hp}.3.60 &= m.c.(45-40)  
			\end{cases}
			\Rightarrow 
			\begin{cases}
				\left(10^3 - \mathcal{P}_{hp} \right).5.60  &= m.4200.25\\
				\mathcal{P}_{hp}.3.60 &= m.4200.5  
			\end{cases} 
			\Rightarrow
			\begin{cases}
				\mathcal{P}_{hp}  &= 250\ W \\
				m &= \dfrac{15}{7}\ kg  
			\end{cases}
			$
			\itemch 
			\itemch $(1000-250)(\tau _3 - 8.60) = \dfrac{15}{7}.4200.(100-40)  
			\Rightarrow \tau _3 = 1200\ s = 20$ phút.
		\end{itemchoice}
	}
\end{ex}
\subsubsection*{PHẦN 3. Thí sinh trả lời từ câu 1 đến câu 6.}
\setcounter{ex}{0}
\begin{ex}% Câu 1
	Một nhà bác học đang tiến hành đo thí nghiệm trong phòng. Trong nhật ký thí nghiệm, nhà bác học ghi như sau: "Lần thí nghiệm thứ $101$, nhiệt độ phòng thí nghiệm đang là $295\ K$". Nhiệt độ phòng lúc đó tương ứng là bao nhiêu độ Celsius (làm tròn kết quả đến chữ số hàng đơn vị)?
	\SA[oly]{22}
	\loigiai{
		\begin{itemize}
			\item $t(^\circ C) = T(K) - 273 = 295 - 273 = 22^\circ C$
		\end{itemize}
	}
\end{ex}

\begin{ex}% Câu 2
	Một động cơ nhiệt có hiệu suất $40\%$, tác nhân đã nhận từ nguồn nóng nhiệt lượng $Q_1 = 1,5\ MJ$, nhiệt lượng truyền cho nguồn lạnh bằng bao nhiêu $MJ$?
	\SA[oly]{0,9}
	\loigiai{
		\begin{itemize}
			\item $A = H.Q_1 = 0,4.1,5 = 0,6\ MJ$
			\item $Q_2 = Q_1 - A = 1,5 - 0,6 = 0,9\ MJ$
		\end{itemize}
	}
\end{ex}

\begin{ex}% Câu 3
	Đổ $100\ g$ nước ở $10^\circ C$ vào một bình chứa ở nhiệt độ $t$. Sau khoảng thời gian đủ dài, bình chứa $5\ g$ nước đá ở $0^\circ C$ và $95\ g$ nước ở $0^\circ C$. Biết bình chứa có nhiệt dung $150\ J/K$, nước có nhiệt dung riêng là $4,2\ J/g.K$, nước đá có nhiệt nóng chảy riêng là $330\ J/g$. Coi không có sự trao đổi nhiệt với môi trường. Giá trị của $t$ là bao nhiêu $^\circ C$?
	\SA[oly]{-39}
	\loigiai{
		\begin{itemize}
			\item Sau khi cân bằng nhiệt có cả nước và nước đá nên nhiệt độ cân bằng là $0^\circ C$.
			\item Phương trình cân bằng nhiệt: $Q_{thu} = Q_\text{tỏa} \Rightarrow C_b.(0-t) = m_n.c_n.\Delta t + m_\text{đ}.\lambda \Rightarrow 150.(0-t) = 100.4,2.10 + 5.330 \Rightarrow t = -39^\circ C$.
		\end{itemize}
	}
\end{ex}

\begin{ex}% Câu 4
	Một nhiệt lượng kế chứa $3$ lít rượu ở $18^\circ C$. Người ta thả một quả cầu đồng có khối lượng $200\ g$ được nung nóng tới $100^\circ C$ vào nhiệt lượng kế. Thả thêm lần lượt từng quả cầu như thế nhưng được nung nóng ở $90^\circ C$ vào nhiệt lượng kế cho đến khi nhiệt độ của rượu đạt trên $25^\circ C$. Bỏ qua sự trao đổi nhiệt của rượu, quả cầu với nhiệt lượng kế và môi trường. Cho nhiệt dung riêng của đồng là $380\ J/kg.K$, nhiệt dung riêng của rượu là $2500\ J/kg.K$, khối lượng riêng của rượu là $D=0,8\ g/cm^3$. Số quả cầu tối thiểu cần phải thả thêm bằng bao nhiêu?
	\SA[oly]{8}
	\loigiai{
		\begin{itemize}
			\item $m_r = V.D = 3.10^{-3}.800 = 2,4\ kg$
			\item Phương trình cân bằng nhiệt $m_r.c_r.\Delta t_1 = m_d.c_d.\Delta t_2 + n.m_d.c_d.\Delta t_3
			\Rightarrow 2,4.2500.(25-18) = 0,2.380.(100-25) + n.0,2.380.(90-25)
			\Rightarrow n \approx 7,3$.
			\item Vậy cần thả thêm tối thiểu 8 quả cầu để nhiệt độ của rượu đạt trên $25^\circ C$.
		\end{itemize}
	}
\end{ex}

\begin{ex}% Câu 5
	Ba cốc chứa nước cùng một loại có nhiệt độ lần lượt là $10^\circ C;\ 30^\circ C;\ 45^\circ C$. Nếu rót một nửa lượng nước trong cốc 2 vào cốc 1 thì nhiệt độ cân bằng là $15^\circ C$. Nếu rót một nửa lượng nước trong cốc 2 vào cốc 3 thì nhiệt độ cân bằng là $35^\circ C$. Bỏ qua nhiệt dung của cốc. Nếu rót một nửa lượng nước trong cốc 1 và một nửa lượng nước trong cốc 2 vào cốc 3 thì nhiệt độ cân bằng là bao nhiêu $^\circ C$? 
	\SA[oly]{22,5}
	\loigiai{
		Chuẩn hóa $m_1=1$.
		\begin{itemize}
			\item 
			$\begin{cases}
				m_1c(15-10)=0,5.m_2c(30-15)\\
				m_3c(45-35)=0,5.m_2c(35-30)
			\end{cases}
			\Rightarrow 
			\begin{cases}
				m_2 = \dfrac{2}{3}\\
				m_3 = 0,25m_2 = \dfrac{1}{6}
		\end{cases}$
		\item Phương trình cân bằng nhiệt: $0,5.m_1.c(10-t) + 0,5.m_2.c(30-t)+m_3.c(45-t) = 0 \Rightarrow t = 22,5^\circ C$.
		\end{itemize}
	}
\end{ex}

\begin{ex}% Câu 6
	Hai bình nhiệt lượng kế giống nhau chứa cùng một lượng chất lỏng $X$ ở cùng nhiệt độ.
	\begin{itemize}
		\item Đổ nước có nhiệt độ bằng nhiệt độ của $X$ vào bình 1 rồi thả một mẩu hợp kim vào bình đó thì mực nước đầy đến miệng bình. Khi cân bằng nhiệt thì nhiệt độ chất lỏng trong bình tăng thêm $4^\circ C$, nhiệt độ mẩu hợp kim giảm $70^\circ C$.
		\item Thả 7 mẩu hợp kim giống như trên vào bình 2 thì mực chất lỏng $X$ cũng đầy bình. Khi cân bằng nhiệt thì độ tăng nhiệt độ của chất lỏng $X$ bằng độ giảm nhiệt độ của 7 mẩu hợp kim.
	\end{itemize}
	Cho biết nhiệt dung riêng của nước $c_0 = 4200\ J/kg.K$, khối lượng riêng của nước $D_0 = 1\ g/cm^3$, của hợp kim $D = 3\ g/cm^3$, của chất lỏng $X$ là $D_X$ với $D>D_X>D_0$. Nhiệt dung riêng của hợp kim bằng bao nhiêu $J/kg.K$?
	\SA[oly]{800}
	\loigiai{
		\begin{itemize}
			\item Gọi nhiệt dung của chất lỏng $X$ trong bình nhiệt lượng kế là $C_X$
			\item Thể tích của lượng đổ vào bình 1 bằng thể tích của 6 mẩu hợp kim $\Rightarrow$ $V_0 = 6V$
			\item Cân bằng nhiệt ở bình 1 được $C_X.\Delta t + m_0.c_0.\Delta t = m.c.\Delta t
			\Rightarrow C_X.4 + V_0.D_0.c_0.4 = V.D.c.70
			\Rightarrow C_X.4 + 6V.1.4200.4 = V.3.c.70
			\Rightarrow$ $C_X = 52,5Vc - 25200V (1)$
			\item Cân bằng nhiệt ở bình 2 được $C_X = 7mc \Rightarrow C_X = 21Vc$ (2)
			\item Từ (1) và (2) $52,5Vc - 25200V = 21Vc
			\Rightarrow c = 800\ J/kg.K$
		\end{itemize}
	}
\end{ex}
